\begin{table}[!h]
\centering\centering
\caption{Race gap at the draft margin \label{tab:draft-margin}}
\centering
\resizebox{\ifdim\width>\linewidth\linewidth\else\width\fi}{!}{
\begin{threeparttable}
\begin{tabular}[t]{lcccc}
\toprule
  & (1) Log pick & (2) Log pick & (3) Log pick & (4) Drafted\\
\midrule
P(Black) & -0.288*** & -0.125 & 0.017 & -0.065*\\
 & (0.089) & (0.080) & (0.053) & (0.035)\\
P(other race) & -0.259* & -0.064 & 0.093 & -0.083\\
 & (0.142) & (0.136) & (0.094) & (0.059)\\
\midrule
Observations & 3051 & 3051 & 3051 & 3593\\
R$^2$ & 0.104 & 0.422 & 0.679 & 0.328\\
Within R$^2$ & 0.046 & 0.385 & 0.659 & 0.289\\
Class $\times$ position group FE & Yes & Yes & Yes & Yes\\
Combine, recruit, college (D) & No & Yes & Yes & Yes\\
Pre-draft grade (F) & No & No & Yes & No\\
Sample & Drafted & Drafted & Drafted & Invitees, NFL id\\
Mean of outcome & 4.554 & 4.554 & 4.554 & 0.732\\
\bottomrule
\multicolumn{5}{l}{\rule{0pt}{1em}* p $<$ 0.1, ** p $<$ 0.05, *** p $<$ 0.01}\\
\end{tabular}
\begin{tablenotes}[para]
\item \textit{Notes:} 
\item This table includes the estimation results of a prospect-level analogue of equation (1), $Y_{icg} = \beta_1 Black_i + \beta_2 Other_i + X_i\pi + \delta_{c,g} + \varepsilon_{icg}$ for prospect $i$ of draft class $c$ in position group $g$, where $\delta_{c,g}$ are draft class $\times$ position group fixed effects. The sample is NFL draft prospects of the 2011-2022 classes (drafted players and combine invitees) with a race prediction; 508 prospects without one, including every prospect without a GSIS id, are excluded. Columns (1)-(3): log overall pick among drafted players (a positive coefficient is a later pick). Column (4): linear probability model of being drafted among combine invitees. Race is attached through the GSIS id, which an undrafted invitee has only if he later joined an NFL team, so this sample conditions on later NFL employment, an outcome of the draft: 494 of 1457 undrafted invitees and 9 of 2639 drafted invitees have no GSIS id, and 494 undrafted and 9 drafted invitees have no race measure. The drafted share is 0.644 among all invitees and 0.732 among invitees with a race measure. Column (4) is therefore not an estimate of the race gap in the probability of being drafted. The pre-draft grade is not used in column (4): it exists only for drafted players, so its missing indicator would determine the outcome. Block D: combine measures (40, vertical, bench, broad jump, cone, shuttle, height, weight) and the 247 rating interacted with position group; age at draft; the final college team's Power-conference, SRS and HBCU indicators; and final-season college production with position-group slopes (passing for QBs, rushing and receiving for RBs, receiving for WRs and TEs, tackles, tackles for loss and sacks for DL and LB, tackles, interceptions and passes defended for DBs, kicking and punting; defensive statistics start in 2016). Block F: the CFBD pre-draft grade, which may itself embed scouts' bias. P(Black) and P(other race) come from the draft-free race prediction, whose prior omits the draft-round bucket (the primary prior conditions on it, and it records the outcome); every column includes fixed effects for that prior's covariates (position at NFL entry, entry era, college type, whether a home county is known), as regression calibration requires. Controls with incomplete coverage are set to zero with a missing indicator (interacted like the control). Standard errors are clustered at the college level (final college team as listed in the draft or combine record; prospects without a college form their own cluster). P(Black) and P(other race) enter together, so the coefficient on P(Black) compares a Black with a white player. Race is predicted, not observed: each person's probability of being non-Hispanic Black combines first name, surname and hometown county (BIFSG) with an NFL-specific prior estimated by EM on predetermined characteristics (players: position at entry, rookie era, draft round, college type, county availability; staff: role, unit and era at first appearance); documented race statements are not used (notes/race-prediction-design.md). Regressions use the probabilities (regression calibration) and control for the prior's covariates; the coefficient on P(Black) is the Black-white gap if the probabilities are calibrated given the regression's controls and names and hometown are unrelated to the outcome given race and the controls. The validation exhibit compares the predictions with published TIDES shares.
\end{tablenotes}
\end{threeparttable}}
\end{table}
